\chapter{Related Work}\label{s:related}

This chapter positions this thesis relative to prior work with similar goals
or methodology, and states explicitly how this work differs.

\textbf{Pegasus~\cite{deelman2015pegasus} and Nextflow~\cite{di2017nextflow}.}
Pegasus and Nextflow are mature, production-grade WMSs that solve the
downstream problem addressed by this thesis --- reliable execution,
scheduling, and fault recovery for a workflow that has already been specified
--- extremely well. Neither system, however, accepts natural-language input:
both require the scientist to author an explicit DAG (Pegasus) or dataflow
script (Nextflow) by hand. This thesis does not attempt to replace either
system's execution model; it targets the layer \emph{before} specification,
translating an intent into a workflow specification for a WMS (Brane) whose
execution model is comparable in spirit to Pegasus and Nextflow, but which
additionally enforces data-locality guarantees at the platform level.

\textbf{Common Workflow Language~\cite{crusoe2022cwl}.} CWL takes a
complementary approach to the portability problem: rather than a single
execution engine, it standardises a vendor-neutral workflow description format
so that a single specification can run on multiple WMS back-ends. This thesis
shares CWL's goal of lowering the cost of producing a valid workflow
specification, but attacks it from the opposite direction: instead of a
portable specification format that a human must still write, it generates a
platform-specific (BraneScript) specification directly from natural language.
The two approaches are not mutually exclusive; a natural-language-to-CWL
translator would face largely the same RAG-grounding and validation challenges
studied here, applied to a different target DSL.

\textbf{LLM-based workflow generation (Balis et al.~\cite{balis2024llm}).}
Balis demonstrated that LLMs can generate Parsl-based Python workflows from
natural-language descriptions with promising accuracy on domain-specific
tasks, establishing --- for a different target DSL --- the same basic premise
this thesis relies on: that an LLM can encode enough workflow-pattern
knowledge to serve as a translator. This thesis differs in three respects:
it targets a DSL (BraneScript) with essentially no public pre-training
presence, which required both RAG and supervised fine-tuning rather than
relying on the base model's prior exposure to the target language; it
evaluates translation quality with an explicit three-tier metric (compile,
execute, output match) rather than a single pass/fail judgement; and it
studies a semantic-caching layer as a complementary reproducibility mechanism,
which is outside the scope of that work.

\textbf{Code generation benchmarks (Codex/HumanEval~\cite{chen2021evaluating}).}
Chen et al.\ established the pass@$k$ methodology for evaluating LLM code
generation against a corpus of general-purpose Python problems with unit
tests. The compile/execute/output-match metrics used in
Chapter~\ref{s:evaluation} are a domain-specific adaptation of the same
underlying idea --- functional correctness rather than surface similarity to a
reference --- specialised to a setting (BraneScript, a DSL absent from
HumanEval-style corpora, with side effects on real datasets) where a
general-purpose code benchmark does not apply.

\textbf{Retrieval-Augmented Generation~\cite{lewis2020rag}.} Lewis et al.\
introduced RAG as a general mechanism for grounding LLM generation in an
external knowledge base at inference time, originally in the context of
open-domain question answering. This thesis applies the same underlying
mechanism to a structurally different retrieval problem: rather than
retrieving free-text passages, it retrieves function-level package
specifications and BraneScript syntax fragments, and combines two
independent retrieval indices (language specification and
package/dataset) rather than a single corpus, motivated directly by the
RQ1/RQ2 finding (Chapter~\ref{s:discussion}) that grounding quality, more than
model scale, determines translation accuracy for a low-resource DSL.

\textbf{Intent-based networking~\cite{bn2016ibn}.} The four-phase
expression--translation--activation--assurance lifecycle originally proposed
for network configuration provides the conceptual vocabulary this thesis
adopts for scientific workflow generation (Chapter~\ref{s:background}). The
key difference is domain: network intents are typically drawn from a bounded,
well-specified configuration space, whereas scientific intents may reference
arbitrary datasets, statistical procedures, and domain terminology, which is
why this thesis required a dedicated retrieval and fine-tuning pipeline rather
than a rule-based translator, as is common in the IBN literature.

\textbf{Workflow provenance formats
(CWLProv~\cite{soiland2018cwlprov}, RO-Crate~\cite{soiland2022rocrate}) and
FAIR workflows~\cite{goble2020fair}.} These works define how to record and
package \emph{retrospective} execution provenance and reusable workflow
artefacts once a workflow has been specified and run. This thesis's semantic
cache (Chapter~\ref{s:design}) is best understood as a step towards
\emph{intent} provenance --- recording which natural-language request produced
which workflow --- which none of these formats natively capture (gap G1,
Chapter~\ref{s:background}); a natural extension of this work would export
cache entries in an RO-Crate-compatible structure so that intent provenance
and execution provenance can be queried together.
